\documentclass[../../main]{subfiles}
\begin{document}

This chapter presents the design choices underlying the initial benchmark, which serves as the basis for refining subsequent experiments and assessing the current capabilities and potential of the proposed solver-selection approach.

The chapter is organized into six sections, each addressing a key design decision. The first section discusses the provider selection. The second describes the selection of large language models used for testing. The third section presents the initial prompt-engineering strategy used to define a consistent prompt format for evaluation. The fourth presents the rationale behind the selection of benchmark problems used to test model performance. The fifth explains the metrics adopted to evaluate the testing results. The sixth describes the experiment setup and the pre-processing methods adopted to make automatic testing possible.



\section{Provider Choice}
The first step was to determine which LLMs could be used for the algorithm-selection task. Given the limited computational resources available during testing, we relied on externally hosted LLMs accessed through usage-based APIs. The selection prioritized generous free tiers, permissive rate limits, and straightforward integration. This led to the choice of the following providers:
\begin{itemize}
    \item \texttt{Gemini API v1}~\cite{GeminiAPI}, offered by Google DeepMind. Gemini is a family of LLMs with multiple sizes and capabilities. This provider was selected for its strong reasoning abilities, robust tool-use features, and overall high-quality text generation. For the purpose of this research, the version v1~\cite{GeminiAPIver} was preferred as it was more stable and the primary concern of this work was its text-generation capability.

    \item \texttt{Groq API}~\cite{GroqAPI}, provided by Groq. Groq offers
    high-performance inference solutions through its specialized hardware
    architecture. The Groq API exposes a selection of LLMs through a simple and
    lightweight interface, enabling fast and low-latency experimentation.
\end{itemize}

Both APIs were selected for their ease of use, flexibility, overall performance, and, critically, their comparatively generous rate limits relative to competing services. The rate limits of both APIs are reported in \Cref{tab:rate_limits_gemini} and \Cref{tab:rate_limits_groq}. 

\begin{table}[ht]
    \centering
    \resizebox{\textwidth}{!}{
    \begin{tabular}{lllll}
        \hline
        Model & RPM & RPD & TPM & TPD \\
        \hline
        \texttt{allam-2-7b} & 30 & 7000 & 6000 & 500000 \\
        \texttt{deepseek-r1-distill-llama-70b} & 30 & 1000 & 6000 & 100000 \\
        \texttt{gemma2-9b-it} & 30 & 14400 & 15000 & 500000 \\
        \texttt{groq/compound} & 30 & 250 & 70000 & -- \\
        \texttt{groq/compound-mini} & 30 & 250 & 70000 & -- \\
        \texttt{llama-3.1-8b-instant} & 30 & 14400 & 6000 & 500000 \\
        \texttt{llama-3.3-70b-versatile} & 30 & 1000 & 12000 & 100000 \\
        \texttt{meta-llama/llama-4-maverick-17b-128e-instruct} & 30 & 1000 & 6000 & 500000 \\
        \texttt{meta-llama/llama-4-scout-17b-16e-instruct} & 30 & 1000 & 30000 & 500000 \\
        \texttt{meta-llama/llama-guard-4-12b} & 30 & 14400 & 15000 & 500000 \\
        \texttt{meta-llama/llama-prompt-guard-2-22m} & 30 & 14400 & 15000 & 500000 \\
        \texttt{meta-llama/llama-prompt-guard-2-86m} & 30 & 14400 & 15000 & 500000 \\
        \texttt{moonshotai/kimi-k2-instruct} & 60 & 1000 & 10000 & 300000 \\
        \texttt{moonshotai/kimi-k2-instruct-0905} & 60 & 1000 & 10000 & 300000 \\
        \texttt{openai/gpt-oss-120b} & 30 & 1000 & 8000 & 200000 \\
        \texttt{openai/gpt-oss-20b} & 30 & 1000 & 8000 & 200000 \\
        \texttt{playai-tts} & 10 & 100 & 1200 & 3600 \\
        \texttt{playai-tts-arabic} & 10 & 100 & 1200 & 3600 \\
        \texttt{qwen/qwen3-32b} & 60 & 1000 & 6000 & 500000 \\
        \hline
    \end{tabular}
    }
    \caption{Rate limits for models available through the Groq API~\cite{GroqRL}. Each row corresponds to a distinct model, identified by its API name. The columns report the enforced usage constraints: requests per minute (RPM), requests per day (RPD), tokens per minute (TPM), and tokens per day (TPD). A dash (--) indicates that no explicit daily token limit is enforced for that model.}
    \label{tab:rate_limits_groq}
\end{table}

\begin{table}[ht]
    \centering
    \small
    \begin{tabular}{@{} >{\raggedright\arraybackslash}p{6cm} r r r r @{}}
        \toprule
        Model & RPM & RPD & TPM & TPD \\
        \midrule
        \texttt{gemini-2.5-pro} & 5 & 100 & 250000 & -- \\
        \texttt{gemini-2.5-flash} & 10 & 250 & 250000 & -- \\
        \texttt{gemini-2.5-flash-lite} & 15 & 1000 & 250000 & -- \\
        \texttt{gemini-2.0-flash} & 15 & 200 & 1000000 & -- \\
        \texttt{gemini-2.0-flash-lite} & 30 & 200 & 1000000 & -- \\
        \bottomrule
    \end{tabular}
    \caption{Rate limits for models available through the Gemini API~\cite{GeminiRL}. All the columns are organized as in \Cref{tab:rate_limits_groq}.}
    \label{tab:rate_limits_gemini}
\end{table}

\section{Large Language Models Selection}

Both providers offer a broad set of LLMs with varying capabilities and constraints, so an initial filtering step was required. Several options were excluded because they are not designed for text generation. In particular, \texttt{playai-tts} and \texttt{playai-tts-arabic} are text-to-speech LLMs, and therefore unsuitable for testing.

Additional models were excluded because they are currently unavailable or deprecated: \texttt{deepseek-r1-distill-llama-70b}, \texttt{gemini-2.0-flash-lite}, and \texttt{gemma2-9b-it}.

Two more LLMs were excluded due to insufficient context window size. Although their rate limits were acceptable, their token capacity was too small to accommodate even a single full MiniZinc model as input: \texttt{meta-llama/llama-prompt-guard-2-22m} and \\\texttt{meta-llama/llama-prompt-guard-2-86m}.

Finally, \texttt{allam-2-7b} was removed because it failed to follow instructions consistently, often producing incomplete, inconsistent, or unreadable outputs.


After this filtering stage, 18 LLMs remained as a stable base for the
evaluation phase.

\section{Prompt General Structure}
\label{sec:promptStruct}
To determine which LLM was best suited for algorithm selection, it was necessary to design a consistent prompt format to query each model.
The primary objective was to define a structure that was as short and clean as possible, for two main reasons:
\begin{itemize}
    \item Minimise prompt-induced bias: An excessively descriptive or overly complex prompt could negatively influence model behaviour, since long prompts may lead to phenomena such as ``context rot''~\cite{ContextRot}: a progressive decay in response quality as prompt length increases.
    \item Reduce token usage: Since the testing setup is constrained by API rate limits, keeping the prompt compact minimises token consumption.
\end{itemize}

\subsection{Output Structure}
Ensuring a standardized output format was equally important.
Automated testing requires that model outputs follow a strict and predictable format. Any deviation introduces ambiguity during parsing and prevents reliable extraction of solver selections. Maintaining this structure is therefore essential to ensure consistent and fully automated evaluation.

Large or verbose responses also impose practical limitations on the available context window. Because each message contributes to the total token count, excessively long outputs reduce the room available for subsequent turns and larger prompts.

For these reasons, the output format was fixed as an array of three strings.
\[
[``1^{st}\text{Solver}'',\ ``2^{nd}\text{Solver}'',\ ``3^{rd}\text{Solver}'']
\]

Selecting the top three solvers enables two forms of evaluation:
\begin{itemize}
    \item Single-solver evaluation: Measures whether the solver chosen by the LLM is the single best solver for the given instance. If it is not, the evaluation can quantify how close its performance is to the optimal solver.
    \item Parallel-solver evaluation: Measures the effectiveness of running the top three solvers selected by the LLM in parallel. The best result among the three is considered, allowing assessment of whether any of them corresponds to the single best solver for the instance, or, if not, how close the best among the three comes to the optimal performance.
\end{itemize}
The metrics used for these evaluations will be detailed in Section~\ref{sec:metrics}.

Taking all of these considerations into account, the resulting prompt structure is as follows:

\begin{definition}{Prompt Structure}
    MiniZinc model:\\
    \ldots \texttt{MiniZinc problem model (\texttt{.mzn} content)} \ldots

    \vspace{0.5em}
    MiniZinc data:\\
    \ldots \texttt{Instance relative data (\texttt{.dzn} or \texttt{.json} content)} \ldots

    \vspace{0.5em}
    The goal is to determine which constraint programming solver would be best suited for this problem, considering the following options:

        - $s_1$,\\
        - $s_2$,\\
        - \ldots\\
        - $s_n$
    where $s_1 \ldots s_n \in SolverList$.

    Answer only with the name of the 3 best solvers inside square brackets separated by comma and nothing else.
\end{definition}


% \begin{figure}[ht]
%     \centering
%     \includegraphics[width=\linewidth]{PromptExample.png}
%     \caption{Example of prompt}
%     \label{fig:promptEx}
% \end{figure}

\section{Problem Selection}
\label{sec:probSelection}
A crucial component of the testing pipeline is the selection of benchmark problems. Consistent and meaningful evaluation requires a set of benchmark problems that are reliable, diverse, and representative of real solver behavior. To meet these requirements, the problem set should satisfy the following criteria:
\begin{itemize}
    \item Extensive prior testing: The problems must be validated and
    associated with reliable solver performance data, preferably obtained from
    recent evaluations of state-of-the-art solvers.
    \item Diversity: The set must include a varied mix of
    problem types, including combinatorial problems, real-world applications,
    and puzzle-like tasks-covering all major categories: Maximization, Minimization and Satisfaction.

    This ensures that LLM performance can be assessed across different solving
    paradigms.
    \item Complexity: The problems must be sufficiently challenging so
    that solver selection is non-trivial and the LLM's reasoning abilities are
    meaningfully tested.
\end{itemize}
Based on these criteria, the selected benchmark was the problem set from the MiniZinc Challenge 2025~\cite{PhilMznChallenge, MznProblems, MznResults25}. These problems are specifically curated to benchmark the strongest solvers of the year and therefore represent an ideal test bed for algorithm selection.

The problem set contains twenty problems: 1 satisfaction problem, 3 maximization problems, 16 minimization problems.

Each problem consists of a \texttt{.mzn} file containing the MiniZinc model~\cite{MiniZinc}, which defines the decision variables, constraints, and objective function. Every problem is further accompanied by five data instances, each stored in either a \texttt{.dzn} or a \texttt{.json} file providing the problem-specific parameters and constants, yielding a total of 100 diverse and challenging test scenarios.

\section{Test Metrics}
\label{sec:metrics}
In order to evaluate LLM performance rigorously, it is necessary to select a standard metric for assessing the quality of each model's output. In addition, a comparative metric is required to evaluate how well each LLM performs relative to the Single Best Solver (SBS).

Before analysing the evaluation metrics, the systems to which these metrics apply must be defined, namely, the solvers. In this context, a solver is a program that takes as input a MiniZinc description of a computational problem together with the associated instance data, and returns an observable outcome consisting of zero or more solutions to the given problem.  

For decision problems, the outcome may simply be ``satisfiable'' or ``unsatisfiable'', while for optimisation problems the outcome of interest is typically the best objective value found within the allotted time. An evaluation metric, or performance metric, is a function that maps the outcome of a solver on a given instance to a numerical value quantifying the solver's performance on that instance. 

An evaluation metric is often not just defined by the output of the solver. Indeed, it can be influenced by other actors, such as the computational resources available, the problems on which we evaluate the solver, and the other solvers involved in the evaluation. For example, it is often unavoidable to impose a \texttt{timeout} \(\tau\) on the solver's execution when termination within a reasonable time cannot be guaranteed (e.g., NP-hard problems). Timeouts make the evaluation feasible but inevitably couple the evaluation metric to the execution context. For this reason, the evaluation of a meta-solver should also consider the scenario that encompasses the solvers to evaluate, the instances used for the validation, and the timeout. Formally, at least for the purposes of this paper, we can define a scenario as a triple (\(\mathcal{I}, \mathcal{S},\tau\)), where: \(\mathcal{I}\) is a set of problem instances, \(\mathcal{S}\) is a set of individual solvers, \(\tau \in (0,+\infty) \) is a timeout such that the outcome of every solver \(s \in \mathcal{S}\) on every instance \(i \in \mathcal{I}\) is always measured within the time interval \([0,\tau]\).

Evaluating meta-solvers over heterogeneous scenarios (\(\mathcal{I}_1, \mathcal{S}_1,\tau_1\)), (\(\mathcal{I}_2, \mathcal{S}_2,\tau_2\)), \dots, is complicated by the fact that the sets of instances \(\mathcal{I}_k\), the sets of solvers \(\mathcal{S}_k\) and the timeouts \(\tau_k\) can be very different. Scenarios that involve optimization problems may even present additional complexities.

To these ends, two separate metrics were adopted.

\subsection{Metric for Solver Score}
\label{sec:scoreMetric}

\begin{figure}[h]
    \centering
    \includegraphics[width=0.5\linewidth]{TableSolverMetricsExample.png}
    \caption{Solver performances example}
    \label{fig:solverperformaceex}
\end{figure}

We can now proceed to associate, for each instance $i$ and solver $s$, a weight that
quantitatively represents how well $s$ performs on instance $i$ within time $T$. We
define the scoring value of $s$ (henceforth, the score) on instance $i$ at a
given time $t$ as a function $\mathrm{score}_{\alpha,\beta}$~\cite{evaluationMetaSolvers, PortfolioAproaches} defined as follows:

\[
\mathrm{score}_{\alpha,\beta}(s,i,t)=
\begin{cases}
0, &
    \text{if }\mathrm{sol}(s,i,t)=\mathrm{unk},
    \\[6pt]

1, &
    \text{if }\mathrm{sol}(s,i,t)\in\{\mathrm{opt},\mathrm{uns}\},
    \\[6pt]

\beta, &
    \begin{aligned}
    &\text{if }\mathrm{sol}(s,i,t)=\mathrm{sat}
    \\[-2pt]
    &\text{ and }\mathrm{MIN}(i)=\mathrm{MAX}(i),
    \end{aligned}
    \\[10pt]

\displaystyle
\max\!\left\{
  0,\;
  \beta-(\beta-\alpha)
  \frac{\mathrm{val}(s,i,t)-\mathrm{MIN}(i)}
       {\mathrm{MAX}(i)-\mathrm{MIN}(i)}
  \right\}, &
    \begin{aligned}
    &\text{if }\mathrm{sol}(s,i,t)=\mathrm{sat}
    \\[-2pt]
    &\text{and $i$ is a minimization problem},
    \end{aligned}
    \\[12pt]

\displaystyle
\max\!\left\{
  0,\;
  \alpha+(\beta-\alpha)
  \frac{\mathrm{val}(s,i,t)-\mathrm{MIN}(i)}
       {\mathrm{MAX}(i)-\mathrm{MIN}(i)}
  \right\}, &
    \begin{aligned}
    &\text{if }\mathrm{sol}(s,i,t)=\mathrm{sat}
    \\[-2pt]
    &\text{and $i$ is a maximization problem}.
    \end{aligned}
\end{cases}
\]


Here, $\mathrm{MIN}(i)$ and $\mathrm{MAX}(i)$ denote the minimal and maximal
objective function values found by any solver $s$ at the time limit $T$.

\medskip

As an example, consider the scenario in \Cref{fig:solverperformaceex} showing three different solvers
on the same minimization problem. Let $T=500$, $\alpha=0.25$, $\beta=0.75$.
Solver $s_1$ finds the optimal value (40), therefore it receives score $0.75$.
Solver $s_2$ finds the maximal value (50), hence score $0.25$. Solver $s_3$
does not find a solution in time, giving score $0$. If instead $T=800$,
the value of $s_1$ becomes $0.375$ and $s_3$ gets $0.75$. If $T=1000$,
since $s_3$ improves the objective to 10 (marked with a star in the figure),
it receives the highest score.

The parameters used for score calculation in the experiments are: $T=1{,}200{,}000$\,ms (i.e., 20 minutes, the time limit adopted in the MiniZinc Challenge for solver evaluation), $\alpha=0.25$, and $\beta=0.75$.

\subsection{Closed Gap}
\label{sec:closedGap}
Once the solver score metric has been defined, a comparative metric is also required to assess the overall performance of a meta-solver relative to the reference baselines.
To this end, the closed-gap (CG)~\cite{evaluationMetaSolvers} measure was adopted as the evaluation metric.
CG is a relative, meta-solver-specific measure that was employed in the 2015 ICON and 2017 OASC~\cite{lindauer2019}
challenges to handle the disparate nature of the scenarios. 

CG assigns a value to a meta-solver in $(-\infty, 1]$ 
proportional to how much it closes the gap between the best individual solver available, 
or single best solver (SBS), and the virtual best solver (VBS), i.e., an 
oracle-like meta-solver always selecting the best individual solver. The closed gap is a ``meta-metric'', defined in terms of a base evaluation metric $m$ to minimise, which in this context is the scoring metric introduced in Section~\ref{sec:scoreMetric}. 
Formally, if $(I, S, \tau)$ is a scenario then 
\[
m(i, \mathrm{VBS}, \tau) = \min \{ m(i, s, \tau) \mid s \in S \}
\quad\text{for each } i \in I,
\]
and 
\[
\mathrm{SBS} = \arg\min_{s \in S} \sum_{i \in I} m(i, s, \tau).
\]

With these definitions, CG can be defined as follows: Let ($\mathcal{I}, S, \tau$) be a scenario and 
\[
m : \mathcal{I} \times \bigl(S \cup \{S, \mathrm{VBS}\}\bigr) \times [0,\tau] \to \mathbb{R}
\]
an evaluation metric to minimize for that scenario, where $S$ is a
meta-solver over the solvers of~$S$.
Let
\[
m_\sigma \;=\; \sum_{i\in\mathcal{I}} m(i,\sigma,\tau)
\qquad\text{for }\sigma\in\{S,\mathrm{SBS},\mathrm{VBS}\}.
\]
The CG of $S$ with respect to $m$ on that scenario is
\[
\frac{m_{\mathrm{SBS}} - m_S}{\,m_{\mathrm{SBS}} - m_{\mathrm{VBS}}\,}.
\]

\noindent

The assumption $m_{\mathrm{VBS}} > m_{\mathrm{SBS}}$ is required, i.e., no single-solver can be the $\mathrm{VBS}$ (otherwise, no algorithm selection would be needed, given that its objective is to reach the $\mathrm{VBS}$). Unlike other scores, CG is designed specifically for meta-solvers. Applying it to individual solvers would assign~0 to the $\mathrm{SBS}$ and a negative score to the remaining solvers,
proportional to their performance difference with respect to the $\mathrm{SBS}$ and
the gap $m_{\mathrm{SBS}} - m_{\mathrm{VBS}}$, which makes little sense for
individual solvers, as it wouldn't reflect their actual performance overall. 

\section{Experiment Setup}
\label{sec:expSetup}
Having defined both the prompt format used to query the LLMs (\Cref{sec:promptStruct}) and the benchmark problem set on which they are evaluated (\Cref{sec:probSelection}), the next step is to evaluate them automatically over the full set of selected instances.
To this end, we designed an automated testing pipeline that parallelizes execution by assigning one thread per LLM. For each model, requests are issued sequentially, with five requests per problem, each containing a MiniZinc model and a single instance. %encoded as shown in \Cref{fig:promptEx}

Despite preliminary prompt engineering and model filtering, several MiniZinc instances, particularly their associated data files, still exceeded the providers' rate limits. Given this limitation, additional mechanisms were required to prevent limit violations while still allowing evaluation over the complete instance set.

\subsection{Script Sanitization}
\label{sec:scriptMan}

The most direct way to address oversized requests was to reduce their length. As the prompt itself was already minimal, this required direct manipulation of the MiniZinc model (\texttt{.mzn}) and data files.

A first step consisted of removing all non-essential elements, such as comments (starting with \texttt{\%}~\cite{MiniZinc}), tabs, and unnecessary whitespace. While this reduced token usage and standardised script formatting, it proved insufficient on its own. The principal contributor to token overflow was the presence of large data arrays, which not only increased message length but could also introduce irrelevant information into the context, thereby ``distracting'' the LLM from the most pertinent content~\cite{LLMdistract}.

To mitigate this issue, data arrays were truncated to a fixed maximum length of 30 elements, with an inline comment indicating the original size:
\[ 
[e_1,e_2,\dots,e_{30} \texttt{// array too long to display, dimensions: (150)} ]
\]
While effective for simple arrays of scalar values, this approach did not account for the internal complexity of individual elements and performed poorly on more structured data. For this reason, a second truncation mechanism was introduced based on raw character count. Arrays exceeding 120 characters were truncated accordingly, using the same inline annotation to preserve information about the original size. 

\subsection{Custom Delays}
Since each experiment involved multiple problems and multiple sequential requests per LLM, rate limits could still be exceeded even when individual requests were within bounds. To address this, custom delays were introduced into the experiment orchestration logic.

When an error message is received, for example:

\begin{lstlisting}
Error code: 413 - Request too large for model `openai/gpt-oss-120b` in organization `org_01k9qqesvte4d9h5jnhmzvbmy4` service tier `on_demand` on tokens per minute (TPM): Limit 8000, Requested 8939, please reduce your message size and try again. Need more tokens? Upgrade to Dev Tier today at https://console.groq.com/settings/billing
\end{lstlisting}

\begin{lstlisting}
Error code: 429 - Rate limit reached for model `openai/gpt-oss-120b` in organization `org_01k9qqesvte4d9h5jnhmzvbmy4` service tier `on_demand` on tokens per day (TPD): Limit 200000, Used 193047, Requested 10632. Please try again in 26m29.328s. Need more tokens? Upgrade to Dev Tier today at https://console.groq.com/settings/billing
\end{lstlisting}

The error code was inspected. Errors $413$ and $429$ indicated that a rate limit had been exceeded. The error message was then parsed to identify the specific limit involved. If the limit concerned tokens per minute (TPM) or requests per minute (RPM), the system paused execution for 60 seconds before retrying. If the exceeded limit was tokens per day (TPD) or requests per day (RPD), the message was further analysed to extract the cooldown duration, expressed in the form $XX$h\,$XX$m\,$XX.XX$s, where $h$ denotes hours, $m$ minutes, and $s$ seconds. The required delay was then computed from this value, after which the request was retried.

\end{document}



